\section{Постановка задач}

В ходе лабораторной работы \textit{по варианту №5} будут выполнены следующие задачи:

\begin{enumerate}
  \item Приблизить математическое ожидание и дисперсию.
  \item Построить вариационный ряд и эмпирическую функцию распределения (график).
  \item Найти медиану.
  \item Построить гистограмму выборки.
\end{enumerate}

Для анализа предоставлена следующая выборка:

\begin{table}[h!]
\centering
\begin{tabular}{|c|c||c|c||c|c||c|c|}
\hline
№ & Значение & № & Значение & № & Значение & № & Значение \\ \hline \hline
1 & 2.89 & 6  & 3.32 & 11 & 3.09 & 16 & 3.30 \\ \hline
2 & 3.14 & 7  & 3.01 & 12 & 2.66 & 17 & 3.07 \\ \hline
3 & 3.12 & 8  & 2.55 & 13 & 3.21 & 18 & 3.42 \\ \hline
4 & 3.34 & 9  & 3.25 & 14 & 2.95 & 19 & 3.05 \\ \hline
5 & 2.89 & 10 & 3.24 & 15 & 3.38 & 20 & 3.03 \\ \hline
\end{tabular}
\caption{Исходные данные выборки}
\end{table}